\chapter{Análisis exploratorio e inferencia clásica}\label{cap:exploracion}

Antes de modelar conviene conocer la respuesta, cuantificar con intervalos de confianza su
centro y su dispersión, y examinar sus relaciones marginales con cada explicativa. Este capítulo
aplica las herramientas de inferencia del curso \parencite[caps.~6--10]{wasserman2004}.

\section{La respuesta: mortalidad por cáncer}

La mortalidad media es de \res{media} muertes por \num{100000} habitantes, con una desviación
típica de \res{dt}. La \cref{tab:respuesta} recoge los intervalos de confianza al 95\,\% para los
parámetros de posición y dispersión, cada uno construido con la cantidad pivotal adecuada:

\begin{itemize}
  \item para la media, $\bar X \pm t_{n-1,\,\alpha/2}\,S/\sqrt n$, válido por el teorema central
  del límite aunque la distribución no sea normal (con $n=\res{n-condados}$);
  \item para la varianza, el pivote $(n-1)S^2/\sigma^2\sim\chi^2_{n-1}$, que sí exige normalidad;
  por eso se acompaña de estimadores robustos;
  \item para la mediana, un intervalo de libre distribución basado en estadísticos de orden
  (cobertura exacta \res{cobertura-mediana}\,\%) y otro bootstrap percentil
  \parencite{efron1993}, que coinciden.
\end{itemize}

\begin{table}[htbp]
\centering
\small
\caption{Estimación puntual y por intervalos (95\,\%) de la mortalidad por cáncer}\label{tab:respuesta}
\begin{tabularx}{\linewidth}{l r l X}
\toprule
Parámetro & Estimación & IC 95\,\% & Método \\
\midrule
Media & \num{178.66} & $[\num{177.68};\ \num{179.65}]$ & pivote $t$ de Student \\
Mediana & \num{178.1} & $[\num{177.0};\ \num{179.1}]$ & estadísticos de orden (libre distribución) \\
Mediana & \num{178.1} & $[\num{177.0};\ \num{179.1}]$ & bootstrap percentil (\num{2000} réplicas) \\
Media recortada al 10\,\% & \num{178.26} & $[\num{177.32};\ \num{179.22}]$ & bootstrap percentil \\
M-estimador de Huber & \num{178.26} & --- & estimador robusto \\
Varianza & \num{770.1} & $[\num{732.9};\ \num{810.3}]$ & pivote $\chi^2_{n-1}$ \\
Desviación típica & \num{27.75} & $[\num{27.07};\ \num{28.47}]$ & pivote $\chi^2_{n-1}$ \\
MEDA & \num{17.0} & --- & mediana de las desviaciones absolutas \\
\bottomrule
\end{tabularx}

\end{table}


La media, la mediana (\res{mediana}), la media recortada y el M-estimador de Huber
\parencite{huber1964} prácticamente coinciden, de modo que los atípicos no desplazan el centro de
la distribución. La asimetría es leve (\res{asimetria}) y la curtosis en exceso moderada
(\res{curtosis}): las colas son algo más pesadas que las de la normal.

\figura{respuesta}{Distribución de la mortalidad por cáncer: histograma con la densidad normal
ajustada y la estimación núcleo, y gráfico cuantil-cuantil normal.}

\subsection{Normalidad}

Todos los contrastes rechazan la normalidad exacta (\cref{tab:normalidad}). Con más de tres mil
observaciones cualquier desviación mínima es detectable, así que el resultado se interpreta
junto al gráfico Q-Q: la forma es casi simétrica y las desviaciones se concentran en las colas.
Para el modelo de regresión la hipótesis relevante no es la normalidad de la respuesta sino la de
los errores, que se examina en el \cref{cap:diagnostico}. El perfil de verosimilitud de Box-Cox
\parencite{boxcox1964} para la respuesta sin explicativas sitúa $\lambda$ en \res{boxcox-lambda}
\res{boxcox-ic}; la transformación se reconsidera, condicionada a las explicativas, en el
\cref{cap:sensibilidad}.

\begin{table}[htbp]
\centering
\small
\caption{Contrastes de normalidad de la mortalidad}\label{tab:normalidad}
\begin{tabularx}{\linewidth}{X r r}
\toprule
Contraste & Estadístico & $p$ \\
\midrule
Shapiro-Wilk & \num{0.9903} & $<\num{0.001}$ \\
Kolmogorov-Smirnov-Lilliefors & \num{0.0280} & \num{0.001} \\
D'Agostino-Pearson K² & \num{127.3633} & $<\num{0.001}$ \\
Jarque-Bera & \num{269.8196} & $<\num{0.001}$ \\
Anderson-Darling & \num{4.0742} & \num{0.010} \\
$\chi^2$ de Pearson (clases equiprobables) ($k = 50$, 47 g.\,l.) & \num{94.2701} & $<\num{0.001}$ \\
Asimetría (contraste de D'Agostino) & \num{0.275} & $<\num{0.001}$ \\
Curtosis de exceso (contraste de Anscombe-Glynn) & \num{1.350} & $<\num{0.001}$ \\
\bottomrule
\end{tabularx}
\par\smallskip{\footnotesize\raggedright Anderson-Darling informa un $p$ interpolado en tablas (acotado entre 0,01 y 0,15). Con $n$ grande, cualquier desviación mínima resulta significativa: los contrastes se leen junto al gráfico Q-Q.\par}
\end{table}


\subsection{Independencia en el orden del fichero}

Una prueba de rachas sobre los signos de las desviaciones respecto de la mediana, en el orden del
fichero, encuentra \res{rachas} rachas frente a \res{rachas-esperadas} esperadas bajo
independencia ($z=\res{rachas-z}$, \res{rachas-p}). El fichero está ordenado por estado, de modo
que los condados vecinos se parecen: es la primera señal de dependencia geográfica, que
condicionará la inferencia del modelo.

\section{Diferencias entre regiones}

La mortalidad media difiere entre regiones censales (\cref{tab:regiones}): el Sur
(\res{media-sur}) supera a todas las demás y el Oeste (\res{media-oeste}) es la más baja. Como las
varianzas difieren (Levene, \res{levene-p}), la comparación global se hace con el ANOVA de Welch
($F=\res{welch-f}$, \res{welch-gl2} g.\,l. en el denominador, \res{welch-p}), confirmado por el
contraste no paramétrico de Kruskal-Wallis ($H=\res{kruskal-h}$, \res{kruskal-p}). La región
explica el \res{pct-region}\,\% de la variabilidad ($\eta^2=\res{eta2-region}$) y el estado, el
\res{pct-estado}\,\%. Las comparaciones por pares (\cref{tab:pares}) muestran que sólo el Medio
Oeste y el Noreste no difieren entre sí.

\begin{table}[htbp]
\centering
\small
\caption{Mortalidad por región censal}\label{tab:regiones}
\begin{tabular}{l r r l r r}
\toprule
Región & $n$ & Media & IC 95\,\% & Desv. típ. & Mediana \\
\midrule
Sur & \num{1383} & \num{189.0} & $[\num{187.5};$\allowbreak\ $\num{190.4}]$ & \num{27.6} & \num{188.2} \\
Medio Oeste & \num{1029} & \num{174.5} & $[\num{173.0};$\allowbreak\ $\num{176.0}]$ & \num{23.9} & \num{173.5} \\
Noreste & \num{217} & \num{172.5} & $[\num{170.6};$\allowbreak\ $\num{174.4}]$ & \num{14.2} & \num{173.3} \\
Oeste & \num{418} & \num{158.0} & $[\num{155.4};$\allowbreak\ $\num{160.7}]$ & \num{27.2} & \num{157.7} \\
\bottomrule
\end{tabular}
\par\smallskip{\footnotesize\raggedright ANOVA: $F = \num{178.0}$ (3 y \num{3043} g.\,l.), $p < \num{0.001}$, $\eta^2 = \num{0.149}$. Levene (mediana): $p < \num{0.001}$. Welch: $F = \num{164.3}$, $p < \num{0.001}$. Kruskal-Wallis: $H = \num{478.3}$, $p < \num{0.001}$.\par}
\end{table}

\begin{table}[htbp]
\centering
\small
\caption{Comparaciones por pares entre regiones}\label{tab:pares}
\begin{tabular}{l r l r r}
\toprule
Par & Diferencia & IC 95\,\% (Tukey) & $p$ Tukey & $p$ Mann-Whitney (Holm) \\
\midrule
Medio Oeste -- Noreste & \num{2.0} & $[\num{-2.9};\ \num{6.9}]$ & \num{0.731} & \num{0.263} \\
Medio Oeste -- Oeste & \num{16.5} & $[\num{12.6};\ \num{20.3}]$ & $<\num{0.001}$ & $<\num{0.001}$ \\
Medio Oeste -- Sur & \num{-14.5} & $[\num{-17.2};\ \num{-11.7}]$ & $<\num{0.001}$ & $<\num{0.001}$ \\
Noreste -- Oeste & \num{14.5} & $[\num{9.0};\ \num{20.0}]$ & $<\num{0.001}$ & $<\num{0.001}$ \\
Noreste -- Sur & \num{-16.4} & $[\num{-21.2};\ \num{-11.6}]$ & $<\num{0.001}$ & $<\num{0.001}$ \\
Oeste -- Sur & \num{-30.9} & $[\num{-34.6};\ \num{-27.2}]$ & $<\num{0.001}$ & $<\num{0.001}$ \\
\bottomrule
\end{tabular}

\end{table}


\figura{regiones}{Mortalidad por región censal: distribución en cada región y media con su
intervalo de confianza al 95\,\%.}

\section{Dos contrastes de dos muestras}

\begin{itemize}
  \item \textbf{Pobreza.} Los condados con al menos un 20\,\% de población bajo el umbral de
  pobreza (\res{pobreza-n-alta}) tienen una mortalidad media de \res{pobreza-media-alta}, frente a
  \res{pobreza-media-baja} en el resto: una diferencia de \res{pobreza-dif} muertes por
  \num{100000} habitantes \res{pobreza-dif-ic} (Welch, $t=\res{pobreza-welch-t}$), con un tamaño
  del efecto grande ($d=\res{pobreza-d}$).
  \item \textbf{Ensayos clínicos.} Los \res{ensayos-n-con} condados con algún ensayo clínico
  tienen una mortalidad \res{ensayos-dif} puntos menor (Welch \res{ensayos-welch-p};
  Mann-Whitney \res{ensayos-mw-p}). Es una diferencia marginal que, como se verá, desaparece al
  ajustar por las características socioeconómicas: los ensayos se concentran en condados urbanos
  y ricos.
\end{itemize}

\section{Asociaciones marginales}

La \cref{fig:correlaciones} y la \cref{tab:correlaciones} (\cref{ap:tablas}) recogen la
correlación de cada explicativa con la mortalidad, con su intervalo de confianza por la
transformación $z$ de Fisher y $p$ corregidos por Holm \parencite{holm1979}; \res{n-corr-sig} de
las \res{n-corr} son significativas. Las más fuertes son la incidencia ($r=\res{r-incidencia}$), el
porcentaje de adultos con título universitario ($r=\res{r-grado25}$), la renta mediana
($r=\res{r-renta}$), la cobertura sólo pública ($r=\res{r-solopublico}$) y la pobreza
($r=\res{r-pobreza}$). El gradiente por renta es monótono: la mortalidad media pasa de
\res{renta-decil1} en el decil más pobre a \res{renta-decil10} en el más rico ($\rho$ de
Spearman $=\res{renta-spearman}$).

\figura{correlaciones}{Correlación de Pearson de cada explicativa con la mortalidad, con su
intervalo de confianza al 95\,\%.}

Estas correlaciones son marginales: muchas explicativas están muy correlacionadas entre sí
(\cref{fig:matriz-correlaciones}), de modo que parte de cada asociación se debe a las demás. El
modelo de regresión múltiple separa esas contribuciones.

\figura{matriz-correlaciones}{Matriz de correlaciones entre las variables candidatas. Los bloques
de correlaciones intensas (educación, renta, pobreza, cobertura sanitaria y empleo) anticipan
problemas de colinealidad.}
